\documentclass[12pt, a4paper, oneside]{article}

% Codificação e idioma
\usepackage[utf8]{inputenc}
\usepackage[T1]{fontenc}
\usepackage[brazil]{babel}

% Margens ABNT: superior 3cm, inferior 2cm, esquerda 3cm, direita 2cm
\usepackage[
  top=3cm,
  bottom=2cm,
  left=3cm,
  right=2cm
]{geometry}

% Espaçamento 1,5 entre linhas (ABNT NBR 14724)
\usepackage{setspace}
\onehalfspacing

% Fontes e tipografia
\usepackage{times}
\usepackage{microtype}

% Recuo de parágrafo ABNT: 1,25cm; primeiro parágrafo também recuado
\usepackage{indentfirst}
\setlength{\parindent}{1.25cm}

% Numeração de página no canto superior direito (ABNT)
\usepackage{fancyhdr}
\pagestyle{fancy}
\fancyhf{}
\fancyhead[R]{\thepage}
\renewcommand{\headrulewidth}{0pt}

% Formatação dos títulos de seção (ABNT: numeração progressiva, sem ponto)
\usepackage{titlesec}
\titleformat{\section}{\normalfont\fontsize{12}{14}\bfseries\MakeUppercase}{\thesection}{1em}{}
\titleformat{\subsection}{\normalfont\fontsize{12}{14}\bfseries}{\thesubsection}{1em}{}
\titleformat{\subsubsection}{\normalfont\fontsize{12}{14}\itshape}{\thesubsubsection}{1em}{}

% URLs e hyperlinks
\usepackage{url}
\usepackage{hyperref}
\hypersetup{
  colorlinks=true,
  linkcolor=black,
  citecolor=black,
  urlcolor=black
}

% Referências ABNT (deve vir após hyperref)
\usepackage[alf]{abntex2cite}

% Título, autores e data
\title{\textbf{Confiabilidade de sistemas de software distribuídos a partir de sinais OpenTelemetry}}
\author{
  Felippe Veloso Marinho \\
  Josoe Santos Queiroz \\
  Maria Eduarda Sampaio \\
  Matheus Amorim Soares Salles \\
  Raphael Lucas Silva Amaral
}
\date{2026}

\begin{document}

\maketitle
\thispagestyle{fancy}

% Resumo em português (ABNT NBR 6028)
\begin{center}
\textbf{RESUMO}
\end{center}
\begin{spacing}{1.0}
\noindent

\vspace{0.5cm}
\noindent\textbf{Palavras-chave:} palavra1; palavra2; palavra3.
\end{spacing}

\vspace{1cm}

\section{Introdução}

Desde a explosão das tecnologias de núvem da década de 2010 a grande maioria dos sistemas de software modernos tem mantido conexão a internet garantindo diversas vantágens da conectividade, ente elas estão o monitoramento e a fácil reparablidade. O resultado da engenhara de software é invisível e em algumas vezes inacessível. Por isso, de aplicações de software é uma das bases para evoluir produtos conectados a internet de forma assertiva e proativa. 

Pela ótica da engenharia de confiablilidade softwares são um desafio interessante, o custo de reparo é baixo, o impacto de uma falha pode ser global e há quantidades astronômicas de dados de comportamento e falhas em sites onde o monitoramento é feito de maneira adequada. 
Há relato

Contudo análises aprofundadas da confiabilidade de softwares ainda são caros devido ao alto custo de conhecimento necessário para que análises personalizadas e mais realistas sejam implementadas em empresas de pequeno e médio porte. 

\section{Revisão de Literatura}
A avaliação da confiabilidade de software é uma área extensivamente pesquisada mas ainda muito dispersa. 
\subsection{Engenharia de confiabilidade de software}
%% copiado
A prática voltada à confiabilidade de sistemas web é conhecida como \textit{Site Reliability Engineering} (SRE) \cite{google_sre} e é prática comum no Brasil desde a metade da última década. 
\subsection{Monitoramento e manutenção de sistemas de software}

%% copiado 
\citeonline{oppenheimer2003} analisaram relatórios de falhas de três grandes serviços de Internet, distinguindo falhas de componentes de falhas de serviço, e identificaram o erro de operador como a maior causa em dois deles.

\section{Materiais e Métodos}

\section{Resultados}


\section{Conclusão}

\section{Declaração de uso de IA}

\bibliography{references}

\end{document}
